\chapter{Reproducibility Report}

\section{Environment}
The detector is trained with Ultralytics 8.4.45 on PyTorch; the RAG service uses
FastAPI, LlamaIndex, LangChain, Ollama, and ChromaDB; the mobile app is built with
Flutter. Dependency definitions and scripts in the project repositories support
deterministic re-runs.

\section{Deployment}
The RAG service is deployed on a RunPod RTX A5000 instance and exposed to the mobile
app through an
ngrok tunnel; the detector ships as a TensorFlow Lite asset bundled with the Flutter
app. No container image is required for the deployed path.

\section{Seed and Split Management}
Because the dataset is built from video frames, splits are performed
\emph{chronologically within each class/video group} (not by random shuffle) to
minimize near-duplicate frame leakage between train, validation, and test. The v3
detector uses a $70/20/10$ split.

\section{Hardware}
Table~\ref{tab:hardware} reports the hardware used for each role.

\begin{table}[htbp]
\centering
\caption{Hardware and runtime by role.}
\label{tab:hardware}
\begin{tabularx}{\textwidth}{@{}l>{\raggedright\arraybackslash}p{0.30\textwidth}>{\raggedright\arraybackslash}X@{}}
\toprule
\textbf{Role} & \textbf{Hardware} & \textbf{Notes}\\
\midrule
Detector training & NVIDIA GTX 1650 Ti (4\,GB) & Ultralytics 8.4.45 / PyTorch; $\sim$34\,h GPU $+$ $\sim$40\,h hand annotation.\\
RAG experiments   & NVIDIA RTX 3050 Laptop (4\,GB), 16\,GB RAM & Ollama (CUDA), Flash Attention enabled.\\
Production serving & RunPod NVIDIA RTX A5000 (24\,GB VRAM) $+$ ngrok & FastAPI RAG
service; 9 vCPU (Intel Xeon Gold 6342 @ 2.80\,GHz), 50\,GB RAM, 30\,GB storage.\\
\bottomrule
\end{tabularx}
\end{table}

\section{Experiment Traceability}
Detector experiments are traceable through the Ultralytics run directories
(\texttt{results.csv}, \texttt{per\_class\_metrics.csv}) for each iteration
(v1/v2/v2@640/v3). The RAG pipeline is evaluated against a fixed set of 20 standardized
bilingual questions (12 English, 8 Arabic) across camera and text modes.

\section{One-Command Rebuild}
The detector is reproduced by running the training scripts
\texttt{01\allowbreak\_split\allowbreak\_dataset.py} through \texttt{06\_eval.py} and exporting
with \texttt{05\_export\allowbreak\_tflite.py}; the knowledge service is rebuilt with
\texttt{index\_dataset.py} and served with \texttt{api.py}.
